\section{At the End of Every Why}

What ought I ultimately to want?

Not release from creaturehood. Not a clean moral account that makes further surrender unnecessary. Not a feeling by which I can verify that God is near, an outcome by which I can prove prayer effective, or suffering that displays spiritual seriousness. Not even union imagined as the extinction of the one who is united.

I ought to want God.

Not because every other good is unreal, but because every true created good is more truly loved when it is loved in the order of its source and end. Not because justice has become unnecessary, but because the One to whom worship is due has disclosed himself as the One who first loves. ``In this is love, not that we loved God but that he loved us and sent his Son.'' ``We love, because he first loved us.''\endnote{1 John 4:7--21, especially vv.~10 and 19: revealed Trinitarian and Paschal truth, never the reversal ``love is God.''}

When this page is closed, the fact from which the series began will not pause. In that instant I will still be receiving existence. The Giver was there before I noticed, is there while I resist, and will be there at the final created moment. Yet revelation says more than metaphysics can reach. The One who holds me in being has entered human anguish, prayed with a human will, laid down a human life, risen in a human body, and gives that Body as food. The Spirit makes the creature's answer possible from within.

Prayer can therefore become less an attempt to summon an absent God than consent to the presence upon which every attempt already depends. Work can become prayer when it is offered in charity. Communion can become hunger for the One received. Unavoidable suffering can be joined to Christ without being called good; removable suffering can be resisted for love of those he loves. Justice is not abandoned. It is caught up into friendship.

No independent ground has appeared beneath creaturehood---not even in the act of love. The will becomes steady not by finding such ground but by ceasing to demand it. Reverence remains; what recedes is alienation. Service is taken into friendship, and the infinite disproportion between Creator and creature becomes no obstacle to communion.

I give God nothing he lacks. The One whom nothing completes gives himself.

At the end of every creaturely why there is no higher law imposed upon God, no need behind his gift, no cause behind his being. There is no mechanism to master and no wage that could equal the gift. There is the Father who sends, the Son who gives himself, and the Spirit through whom the creature can finally love.

\vspace{1.2\baselineskip}
\begin{center}
{\Large\bfseries God is love.}
\end{center}
